% Shared preamble for the M2 layout prototypes (xelatex + Cardo)
% ⚠ KDP flagged insufficient gutter at 625pp (needs >=0.75in clear from the trim
% edge). The true cause: \plnumber's \llap ruler (below) pokes INTO the inner
% margin by design, so the nominal 0.80in inner left only ~0.66in actually clear
% -- verified by measuring page 21's rendered ink, not by guessing. Fixed by
% moving 0.15in from outer to inner, which leaves \textwidth (inner+outer
% unchanged) untouched -- no reflow, no page-count shift, no cover rebuild.
\usepackage[paperwidth=6in,paperheight=9in,
  inner=0.95in,outer=0.50in,top=0.80in,bottom=0.85in]{geometry}
\usepackage{fontspec}
\setmainfont{Cardo}[
  BoldFont={* Bold}, ItalicFont={* Italic},
  Ligatures=TeX]
\usepackage{graphicx}
\usepackage{xcolor}
\definecolor{rulegray}{gray}{0.55}

% a sense-line: #1 = indent in em, #2 = text; long lines wrap with hanging indent
\newlength{\plhang}\setlength{\plhang}{5em}

% ---- sense-line numbering, which is what ANCHORS both foot-bands ----------
% The 1853's own apparatus keys every reading as PAGE. LINE ("34. 27"), and our
% transcripts are line-faithful, so OUR SENSE-LINE COUNT IS ITS LINE COUNT.
% Verified at printed 34: Τὰ ἔργα τῶν χειρῶν σου μὴ παρίδῃς is line 27, exactly
% as class A says. The apparatus was already speaking this coordinate; the page
% simply was not showing it.
%
% ⚠ ONE NUMBERING SERVES ALL THREE COLUMNS. Part I's Latin recto is line-for-line
% with its Greek verso, and our English is line-keyed to the Greek — verified at
% printed 34, where Greek 27 and English 27 are the same line and both pages run
% to 30. So the number set beside the Greek reads across the whole opening, and
% the Latin is deliberately NOT numbered: a second set of figures on the same
% leaf would say the same thing twice.
%
% ⚠ Numbering counts \pl calls, i.e. SENSE-lines, not typeset lines. That is the
% correct unit and it is also the only stable one — 839 of Part I's lines turn in
% the two columns, and a turned line must not take a number of its own.
\newcounter{plline}\newcounter{plmod}
\newif\ifplnum \plnumfalse
\newcommand{\plreset}{\setcounter{plline}{0}\setcounter{plmod}{0}}

% ⚠ EVERY FIVE LINES IS NOT AN ANCHOR — it is a RULER. A note keyed to line 27 is
% unfindable if the margin only prints 25 and 30: the reader must count, and
% nothing on line 27 says a note is there to count towards. Wilson, on the second
% sample: "there is still no visible anchor."
%
% So there are TWO devices, and they are deliberately different alphabets:
%   · the MARGIN carries ARABIC line figures — the ruler, and the 1853's own
%     PAGE. LINE coordinate, which the apparatus has always cited;
%   · the LINE-END carries a ROMAN superscript — the anchor, which says "there is
%     a note on me" at the place the eye already is.
% Roman against arabic so the two can never be read as the same series. The band
% below repeats the roman, so a marker is matched by looking, not by counting.
%
% A line that carries a note prints its arabic figure too, whatever the every-five
% rule says — the ruler should not go silent exactly where it is wanted.
% ⚠ THE LATIN COLUMN IS NUMBERED ON ITS OWN COUNT, and this was got wrong first
% time. The Latin was left unnumbered on the ground that it is line-for-line with
% the Greek so a second set of figures would repeat them. tools/check_lineparity.py
% shows that is true of 111 of Part I's 131 openings and FALSE of 20 — the plate
% turns its own long lines and the Latin is the wordier column, so it can run five
% or six lines longer than its Greek (printed 48/49, 62/63, 142/143). One set of
% figures would therefore have named the wrong Latin line on one opening in seven,
% and nothing on the page would have shown it.
%
% The MARKER SERIES still runs continuously across the leaf: the Greek's noted
% lines take i…k and the Latin's take k+1…m, so a roman is unique on the page and
% the band needs no column label. \setnotedcol carries the offset.
\newcommand{\notedlines}{}
\newcounter{plofs}
\newcommand{\setnoted}[1]{\renewcommand{\notedlines}{#1}\setcounter{plofs}{0}}
\newcommand{\setnotedcol}[2]{\renewcommand{\notedlines}{#1}\setcounter{plofs}{#2}}
\newif\ifplnoted
\newcounter{plpos}\newcounter{plfound}
\makeatletter
% Walks the noted-line list and records BOTH that this line is noted and its
% position in the list, which is the roman numeral the band will repeat.
\newcommand{\plcheck}[1]{\plnotedfalse
  \setcounter{plpos}{0}\setcounter{plfound}{0}%
  \def\pl@empty{}\ifx\notedlines\pl@empty\else
    \@for\pl@x:=\notedlines\do{\stepcounter{plpos}%
      \ifnum\pl@x=#1\relax\plnotedtrue\setcounter{plfound}{\value{plpos}}\fi}%
  \fi}
\makeatother

% ⚠ NOT \numexpr modulo: \numexpr divides with ROUNDING, so n/5*5=n is true for
% n=3 as well as n=5. Count down in a second counter instead.
\newcommand{\plnumber}{%
  \ifplnum
    \stepcounter{plline}\stepcounter{plmod}%
    \plcheck{\value{plline}}%
    % ⚠ EVERY FIVE ONLY. An earlier pass also printed a figure on every noted
    % line, which was right while the arabic WAS the anchor and became clutter
    % the moment the roman marker took that job — printed 34 came out with
    % figures on eighteen of thirty lines. The two devices divide cleanly: roman
    % anchors, arabic measures.
    \ifnum\value{plmod}=5
      \setcounter{plmod}{0}%
      \llap{{\tiny\color{rulegray}\theplline}\hspace{0.6em}}%
    \fi
  \fi}

% The roman marker, set at the END of the line — \plcheck has already run for this
% line inside \plnumber, so \ifplnoted and \plfound still hold its values.
\newcommand{\plmark}{%
  \ifplnum\ifplnoted
    \textsuperscript{\normalfont
      \romannumeral\numexpr\value{plfound}+\value{plofs}\relax}%
  \fi\fi}
% The number laps into the margin BEFORE the line's own indent, so it sits at the
% block's left edge whatever the indent level — otherwise a deeply indented line
% would carry its figure into the middle of the column.
\newcommand{\pl}[2]{\par\noindent\hangindent=\plhang\hangafter=1
  \plnumber\hspace*{#1em}#2\plmark}
% vertical breathing space between prayer paragraphs
\newcommand{\parasep}{\par\vspace{0.55\baselineskip}}
% small-caps language label used by the stacked layout
\newcommand{\langhead}[1]{\par\vspace{0.4\baselineskip}%
  {\centering\scshape\color{rulegray}#1\par}\vspace{0.35\baselineskip}}
% unit title
\newcommand{\unittitle}[1]{\par\vspace{1.2\baselineskip}%
  {\centering\large\scshape #1\par}\vspace{0.8\baselineskip}}
% thin separator rule
\newcommand{\thinrule}{\par\vspace{0.5\baselineskip}%
  {\centering\color{rulegray}\rule{0.35\textwidth}{0.4pt}\par}\vspace{0.5\baselineskip}}
\raggedbottom
\setlength{\parindent}{0pt}
% bidi must be loaded LAST (provides \RL for the Hebrew headings)
\usepackage{bidi}
